\section{Introduction}


The Koebe--Andreev--Thurston theorem states that any skeleton of triangulation for an oriented compact surface arises as the contact graph of a circle packing on a geometrized surface (\cite[Section~13.7]{Thurston}; see also~\cite{ColindeVerdiere}). In the case of a closed Riemann surface of genus greater than one, one obtains a hyperbolic metric which realizes given combinatorial structure by a circle pattern on the surface. The hyperbolic metric obtained is unique up to isometries, and this theorem is also called \emph{discrete uniformization theorem}~\cite[Chapter~4]{Stephenson}. In this article, we offer yet another way of endowing a hyperbolic metric on a surface in a manner that given weighted graph is realized as the image of a harmonic map with the least energy among a fixed homotopy class as well as all the hyperbolic metrics.



Let $X=(V, E, m_E)$ be a finite weighted graph with $m_E: E \to (0, \infty)$. We identify each edge $e$ with the unit interval $[0, 1]$, and denote by $\wbar e$ its reversed edge. We understand $E$ the set of all oriented edges $e$, $\wbar e \in E$, and the weight function $m_E$ is symmetric, i.e., $m_E(e)=m_E(\wbar e)$.


For any closed Riemann surface $(S, G)$ equipped with a hyperbolic metric $G$, let $f:X \to (S, G)$ be a piecewise smooth map, i.e., the restriction $f_e:[0, 1] \to (S, G)$ is piecewise smooth for each edge $e$. We define the \emph{energy} of $f$ by
\begin{equation}\label{Eq:Dirichlet}
E_G(f):=\frac{1}{2}\sum_{e\,\in\,E}m_E(e)\int_0^1\left\|\frac{d f_e}{dt}(t)\right\|_G^2\,dt.
\end{equation}
It is known that the energy $E_G(f)$ attains the minimum by a harmonic map $h_G$ among the set of all piecewise smooth maps homotopic to $f$ (Section~\ref{Sec:harmonic}). A \emph{harmonic map} $h$ from $X$ into a hyperbolic surface $(S, G)$ is a map such that $h_e$ is a (constant speed) geodesic for each $e \in E$, and satisfies the \emph{balanced condition}
\[
\sum_{e\,\in\,E_x}m_E(e)\frac{d h_e}{dt}(0)=0, \quad \text{for each }x \in V,
\]
where $E_x=\{e \in E \ : \ e_0=x\}$, the set of edges emanating from $x$. Moreover, if a harmonic map $h_G$ is not homotopic to a point nor to a closed curve in $(S, G)$, then $h_G$ is unique as a map due to the negative curvature of $G$. Then, regarding the energy $E_G(h_G)$ of the harmonic map $h_G$ as a function of hyperbolic metrics $G$ on $S$, we may further minimize $E_G(h_G)$ given a fixed homotopy class $\Cc$ of $f$. A central question we consider is whether the energy $E_G(f)$ attains the minimum for a pair $(h_0, G_0)$ of a map $h_0: X \to S$ in $\Cc$ and a hyperbolic metric $G_0$, or not. We show that under some necessary topological condition on $f:X \to S$, the joint minimum of $E_G(f)$ does exist; namely, we solve the double minimization problem where the first minimization is the classical calculus of variation and the second minimization is the variation among the hyperbolic metrics. Furthermore, we identify the hyperbolic metric $G$ when the underlying weighted graph $X$ admits an automorphism group which is compatible with the mapping class group of $S$.

Note that for $i=1, 2$, any two pairs $(h_i, G_i)$ of a hyperbolic metric $G_i$ on $S$ and a harmonic map $h_i:X \to (S, G_i)$ in $\Cc$ have the same energy $E_{G_1}(h_{1})=E_{G_2}(h_{2})$ if there exists an isometry $\f: (S, G_1) \to (S, G_2)$ homotopic to the identity map on $S$ and $h_2=\f \circ h_1$. It holds that if a pair $(h, G)$ attains the minimum of the energy $E_G(h)$, then the hyperbolic metric $G$ on $S$ is unique up to isometries homotopic to the identity map on $S$ under some condition on the fixed homotopy class $\Cc$ as we see below.

We say that a continuous map $f:X \to S$ \emph{fills} $S$ if it is injective and the complement of the image $f$ is a disjoint union of (topological) disks. For example, any skeleton of triangulation of $S$ gives rise to such $X$ and $f$. Let us denote the set of all piecewise smooth maps homotopic to $f$ by $\Cc=[f]$, and the space of hyperbolic metrics on $S$ by $\Mc_{-1}(S)$. We define
\[
\Ec: \Cc \times \Mc_{-1}(S) \to [0, \infty), \qquad (f, G) \mapsto E_G(f).
\]
Note that for any piecewise smooth map $f: X \to (S, G)$, we have $E_G(f)<\infty$, and any continuous map $f$ is homotopic to a piecewise smooth map.

\begin{theo}\label{Thm:main}
Let $X=(V, E, m_E)$ be a connected finite weighted graph, and $S$ be a closed Riemann surface of genus greater than one. If $f:X \to S$ is a continuous map which fills $S$, then there exists a pair $(h_0, G_0)$ which attains the minimum of $\Ec$ on $\Cc \times \Mc_{-1}(S)$, where $\Cc$ is the homotopy class of $f$. Moreover, the hyperbolic metric $G_0$ on $S$ is unique up to isometries homotopic to the identity map on $S$.
\end{theo}



In fact, we prove a more general theorem when $f$ induces a surjective homomorphism $f_\ast: \pi_1(X, x_0) \to \pi_1(S, f(x_0))$ in Theorem~\ref{Thm:general_main}. The unique hyperbolic metric $G_0$ in Theorem~\ref{Thm:main} realizes the least energy harmonic embedding of $X$ into the hyperbolic surface $(S, G_0)$ determined by the weighted graph $X$ and the fixed homotopy class of maps $f: X \to S$.

The problem of minimizing energy~\eqref{Eq:Dirichlet} was studied by Colin de Verdière~\cite{ColindeVerdiereComment} for the embedding of finite weighted graphs into surfaces. The problem of minimizing energy also in hyperbolic metrics was suggested by Kotani and Sunada~\cite[p.~7, Section~2]{KotaniSunadaStandard} and has been reiterated by Sunada~\cite[p.~124, Section~7.7]{SunadaTopological}. They showed the corresponding result in the case of flat tori (of dimension at least $2$); for a connected finite weighted graph $X$, and for an $n$-dimensional torus $\T^n$, if $f:X \to \T^n$ induces a surjective homomorphism from $\pi_1(X)$ to $\pi_1(\T^n)$, then there exists a pair of flat metric $G_0$ on $\T^n$ and a harmonic map $h_0: X \to (\T^n, G_0)$ homotopic to $f$ such that the corresponding energy functional $\Ec$ attains the minimum on $\Cc \times \Mc_0(\T^n)$, where $\Cc=[f]$ and $\Mc_0(\T^n)$ is the space of flat metrics on $\T^n$ whose volume is normalized to $1$ (and such a pair is unique up to translations on $\T^n$). In fact, they proved the result by constructing a flat metric explicitly by using the $1$-homology group of graph $X$, and called the resulting harmonic map $h_0$ into $(\T^n, G_0)$ (or, its equivariant lift on an abelian covering of $X$ into $\R^n$) the \emph{standard realization} of $X$. Then, they proposed the problem in the case of closed hyperbolic surfaces whether such an energy functional has a minimum as a non-Euclidean analogue to their result in dimension~$2$. It seems that a direct construction of such a pair of a metric and harmonic map into closed hyperbolic surfaces would not work by a non-linear nature of target spaces.


Our approach to show the existence of a minimum pair of energy functional $\Ec$ for surfaces is that we define a function $\Ec_\Cc$ associated to $\Ec$ on the Teichm\"uller space $\Tc(S)$. Namely, since the energy $E_G(f)$ attains the minimum at a harmonic map $h_G$ in the homotopy class $\Cc$ of $f$, the function $G \mapsto E_G(h_G)$ is defined on the space of hyperbolic metrics $\Mc_{-1}(S)$. This function naturally yields a function on the Teichm\"uller space $\Tc(S)$ of $S$. Recalling that the Teichm\"uller space $\Tc(S)$ consists of isotopy classes of pairs $((\SS, G_{\SS}), \f)$ of a hyperbolic surface $(\SS, G_{\SS})$ and a diffeomorphism (a marking) $\f:S \to \SS$, we write $G=\f^\ast G_\SS$ the pull-back metric of $G_\SS$ on $S$. Then, the function
\[
\Ec_\Cc([G]):=E_{G}(h_G),
\]
is well-defined on $\Tc(S)$ (Section~\ref{Sec:energy}). We show that the function $\Ec_\Cc$ on $\Tc(S)$ is strictly convex with respect to the Weil--Petersson metric (Theorem~\ref{Thm:convex}) and proper if $f$ induces a surjective homomorphism from $\pi_1(X)$ to $\pi_1(S)$ (Theorem~\ref{Thm:proper}). This shows that a minimum of the function $\Ec_\Cc$ on $\Tc(S)$ exists uniquely and we show that it also yields a minimum pair of the original energy functional $\Ec$.



